\section{Проектирование SLP-векторизатора}
\label{sec:Chapter5} \index{Chapter5}

В данной главе описывается архитектура спроектированного SLP-векторизатора, предназначенного для работы в составе компилятора Go. 
Изложение построено от общих требований и ограничений к деталям каждого этапа обработки.

\subsection{Требования и ограничения для Go-векторизатора}
\label{sec:requirements}

Проектируемый векторизатор должен удовлетворять следующим требованиям, вытекающим из анализа компилятора Go,
проведённого в главе~\ref{sec:Chapter4}.

\begin{itemize}
    \item \textbf{Корректность при явных зависимостях по памяти.} В Go SSA операции с памятью связаны глобальной цепочкой состояний \texttt{Memory}.
    Любое преобразование обязано сохранять этот порядок. Векторизатор должен гарантировать, что объединение нескольких скалярных инструкций в одну векторную не нарушит модель памяти. 
    Консервативный подход требует, чтобы упаковываемые инструкции либо разделяли общий входной \texttt{Memory}-аргумент, либо (для цепочек сохранений) были связаны последовательной зависимостью по памяти.
    
    \item \textbf{Поддержка векторных инструкций ARM64 NEON.} Основной целевой архитектурой на текущем этапе является ARM64 
    с расширением NEON, предоставляющим 128-битные векторные регистры. 
    Векторизатор должен генерировать операции над 128-битными векторами, составленными из двух 64-битных целочисленных элементов.
    Поддержка x86-64 SSE/AVX2 и других архитектур может быть добавлена позже без изменения алгоритмического ядра.

    \item \textbf{Ограничение на упаковку, распаковку и перестановки линий.}
    При отсутствии модели стоимости невозможно оценить, окупит ли выигрыш от векторизации 
    затраты на вспомогательные операции перекомпоновки данных.
    Векторизатор придерживается консервативной стратегии: он не вставляет (или делает лишь очень ограниченное количество) 
    дорогие инструкции упаковки, распаковки и перестановок элементов.

    \item \textbf{Ограничение времени компиляции.} 
    Компилятор Go ориентирован на быструю сборку, поэтому векторизатор не должен добавлять существенных задержек.
    
    \item \textbf{Обработка типов, не содержащих указателей.} Из-за требований точного сборщика мусора 
    векторизации не подлежат указательные типы.
\end{itemize}

\subsection{Общая архитектура}
\label{sec:architecture}

Векторизатор оформлен как проход, работающий на уровне целой функции. 
Он последовательно обходит базовые блоки функции и для каждого блока пытается выделить и
векторизовать группы изоморфных независимых инструкций.

Ключевыми структурами данных являются:
\begin{itemize}
    \item \texttt{pack} --- упорядоченный набор изоморфных значений. Порядок внутри \texttt{pack}'а соответствует линиям будущего векторного регистра.
    Например, \texttt{pack\{v1, v2\}} объединяет две операции, которые станут элементами 0 и 1 векторного регистра.
    Далее для обозначения \texttt{pack} будет использоваться термин --- группа.
    \item \texttt{packSet} --- набор групп, накапливаемый в процессе анализа одного базового блока.
\end{itemize}

Общую схему работы для каждого базового блока можно разделить на три последовательных этапа:
\begin{enumerate}
    \item \textbf{Анализ} -- формирование начальных пар и итеративное наращивание групп потенциально векторизуемых операций с последующим их комбинированием.
    \item \textbf{Валидация} -- разбиение на независимые компоненты и проверка каждой из них на корректность.
    \item \textbf{Генерация кода} -- планирование, легализация и создание векторных инструкций взамен скалярных.
\end{enumerate}

\begin{figure}[htbp]
\centering
\begin{tikzpicture}[
    block/.style={
        rectangle,
        draw,
        rounded corners,
        align=center,
        fill=blue!8,
        draw=blue!40,
        minimum width=4cm,
        minimum height=0.9cm
    },
    container/.style={
        rounded corners,
        draw=black!40,
        dashed, 
        minimum width=11cm,
        inner sep=10pt,
        fill=none
    },
    arrow/.style={->, thick}
]

% Nodes
\node (n1) [block, fill=blue!15, ] {Поиск начальных пар};
\node (n3) [block, fill=blue!15, below=of n1] {Расширение};
\node (n4) [block, fill=blue!15, below=of n3] {Комбинирование};

\node (n5) [block, fill=yellow!20, below=of n4] {Разбиение на компоненты};
\node (n6) [block, fill=yellow!20, below=of n5] {Валидация};

\node (n7) [block, fill=green!15, below=of n6] {Планирование};
\node (n8) [block, fill=green!15, below=of n7] {Легализация};
\node (n9) [block, fill=green!15, below=of n8] {Генерация};

% arrows
\foreach \i/\j in {1/3,3/4,4/5,5/6,6/7,7/8,8/9} {
    \draw [arrow] (n\i) -- (n\j);
}

% Containers (background groups)
\begin{scope}[on background layer]

\node (analysisbox)   [container, fit=(n1)(n3)(n4)] {};
\node (validationbox) [container, fit=(n5)(n6),   ] {};
\node (codebox)       [container, fit=(n7)(n8)(n9)] {};

\end{scope}

% Corner labels (placed inside containers at top‑left)
\node[anchor=north west, font=\small] at ([xshift=4pt, yshift=-4pt]analysisbox.north west) {Анализ};
\node[anchor=north west, font=\small] at ([xshift=4pt, yshift=-4pt]validationbox.north west) {Валидация};
\node[anchor=north west, font=\small] at ([xshift=4pt, yshift=-4pt]codebox.north west) {Генерация кода};

\end{tikzpicture}
\caption{SLP-векторизатор: высокоуровневая схема}
\label{fig:slp_pipeline}
\end{figure}

Ниже подробно описаны шаги каждого этапа.

1. Анализ:
\begin{itemize}
    \item \textbf{Поиск начальных пар.} Сканируются инструкции блока, отыскиваются пары смежных обращений к памяти, удовлетворяющих условиям изоморфизма, независимости и смежности адресов. 
    На их основе формируются начальные пары.
    
    \item \textbf{Расширение групп вверх и вниз.} От начальных пар итеративно достраиваются пары операций для операций-аргументов и операций-потребителей. 
    Процесс продолжается до тех пор, пока удаётся добавлять новые пары. На этом этапе образуется граф групп, отражающий потенциально векторизуемое дерево выражений.
    
    \item \textbf{Комбинирование.} Соседние группы, конец одной из которых совпадает с началом другой, сливаются в более широкие операции. 
    Например, \texttt{pack\{v1, v2\}} и \texttt{pack\{v2, v3\}} превращаются в \texttt{pack\{v1, v2, v3\}}. 
    Это позволяет формировать векторы, ширина которых превышает два элемента, если исходный код содержит большие последовательности однотипных операций.
\end{itemize}

2. Валидация:
\begin{itemize}
    \item \textbf{Разбиение на компоненты связности.} Группы, связанные через аргументы, объединяются в независимые компоненты (подграфы). 
    В дальнейшем каждая компонента обрабатывается отдельно. 
    Это позволяет не отбрасывать сразу все компоненты, если некоторые из них не пройдут дальнейшую проверку.
    
    \item \textbf{Валидация.} Для каждой выделенной компоненты проверяется ряд условий: отсутствие недопустимых внешних использований, замкнутость операндов и когерентность линий. 
    Непрошедшие проверку компоненты отбрасываются.
\end{itemize}

3. Генерация кода:
\begin{itemize}
    \item \textbf{Планирование.} Для прошедшей валидацию компоненты операции упорядочиваются топологически так, 
    чтобы ни одна группа не находилась в выходном наборе раньше тех, от которых она зависит по данным.
    Данный этап упрощает дальнейшую генерацию.
    
    \item \textbf{Легализация.} Группы, длина которых превышает аппаратную ширину векторного регистра (2 элемента для 128-битных векторов при 64-битном скалярном типе), разбиваются на подгруппы.
    
    \item \textbf{Генерация кода.} Выполняется эмиссия векторных SSA-инструкций: для каждой группы создаётся одна векторная инструкция, заменяющая собой исходные скалярные. 
    Внутренние связи между группами заменяются прямыми ссылками на результаты векторных инструкций. 
    Для значений, используемых вне векторизованного фрагмента, вставляются инструкции извлечения отдельных элементов.
\end{itemize}

Далее каждый этап рассматривается подробно.

\subsection{Поиск кандидатов для векторизации}
\label{sec:initial_pairs}

Начальный этап нацелен на выделение пар инструкций, которые послужат основой для дальнейшего расширения. В качестве начальных инструкций служат обращения к памяти --- загрузки и сохранения. 
Этот выбор мотивирован тем, что смежные обращения к памяти предоставляют естественный параллелизм.

Для пары инструкций \((v_1, v_2)\) проверяется следующий набор условий:

\begin{itemize}
    \item \textbf{Изоморфизм}: \(v_1\) и \(v_2\) имеют одинаковый код операции и совместимый тип результата.
    \item \textbf{Векторизуемость операции}: операция входит в «белый список» разрешённых для векторизации операций.
    \item \textbf{Смежность для обращений к памяти}: для инструкций обращения к памяти их адреса должны быть смежными и упорядоченными. Проверка выполняется процедурой, которая использует инфраструктуру прохода \texttt{memcombine}.
    \item \textbf{Отсутствие конфликтов с уже существующими группами}: ни \(v_1\), ни \(v_2\) не должны входить в уже сформированные группы.
    \item \textbf{Независимость}: проверка, что \(v_1\) и \(v_2\) не зависят друг от друга по данным и разделяют общее состояние памяти. 
    Только операции сохранения являются исключением при проверке общего состояния памяти, так как они обязательно будут иметь разные состояния. 
    Для них проверяется, образуют ли они последовательную \texttt{Memory}-цепочку.
\end{itemize}

Дополнительно исключаются инструкции с побочными эффектами. После нахождения пары создаётся новая группа размера 2, которая добавляется в общий набор.

\subsection{Построение графа зависимостей и расширение групп}
\label{sec:expansion}

После формирования начальных пар алгоритм итеративно расширяет множество групп, двигаясь по потоку данных. 
Процесс выполняется до сходимости: за одну итерацию последовательно вызываются процедуры расширения вниз и вверх для каждой группы. 
Если в результате появляются новые группы, итерация повторяется.

\paragraph{Расширение вниз.}
Для группы \(p = \{v_1, v_2\}\) анализируются все инструкции-потребители результатов \(v_1\) и \(v_2\) в пределах того же базового блока. 
Для каждой пары потребителей \((t_1, t_2)\) проверяются условия упаковки. Выбранная пара добавляется в набор как новая группа.

\paragraph{Расширение вверх.}
Для того же \(p = \{v_1, v_2\}\) перебираются все аргументы \(v_1\) и \(v_2\). Применяются те же критерии изоморфизма и независимости.

После завершения расширения накопленные группы могут быть объединены в более длинные. Процедура совмещения просматривает все пары групп: если конец одной совпадает 
с началом другой (например, $p_1 = \{v_1, v_2\}$ и $p_2 = \{v_2, v_3\}$), они сливаются в одну группу (\(p = \{v_1, v_2, v_3\}\)). 
Процедура повторяется до тех пор, пока находятся новые объединения.

Следующим шагом набор групп разбивается на компоненты связности по графу, в котором вершины --- это группы, а ребро существует, 
если одна группа использует в качестве аргумента значение из другой группы. Каждая компонента затем обрабатывается независимо, 
что позволяет векторизовать в одном блоке несколько не связанных друг с другом групп.

\subsection{Валидация}
\label{sec:validation}

Перед генерацией кода каждая выделенная компонента должна пройти ряд проверок корректности.

\subsubsection{Проверка внешних использований}
Для каждого значения, входящего в одну из групп компоненты, анализируются все его использования в функции.
Использование считается допустимым, если оно принадлежит другой группе.

Дополнительно применяется простая эвристика частичной скаляризации. 
Если компонента содержит не менее трёх групп, допускается до несколько значений с внешними пользователями. 
В этом случае предполагается последующая распаковка соответствующего векторного значения для внешнего использования. 
Если лимит исчерпан, компонента отбрасывается.

\subsubsection{Проверка замкнутости операндов}
Для каждого значения проверяются все его операнды. 
Все операнды обязаны принадлежать какой-либо группе той же компоненты. 
Таким образом гарантируется, что вычислительные зависимости не выходят за пределы векторизуемого подграфа.

\subsubsection{Проверка когерентности линий}
Проверяется, что все группы компоненты имеют одинаковую ширину. 
Кроме того, каждое значение должно использоваться в той же линии, в которой оно было создано. 
Для каждого операнда, принадлежащего компоненте, номер его линии обязан совпадать с номером линии инструкции-потребителя. 
Такое ограничение исключает необходимость перестановок (\texttt{shuffle}) между линиями, которые в текущей реализации не поддерживаются.

\subsubsection{Проверка согласованности исходных операндов}
Для каждой группы строится отображение, связывающее позиции аргументов с группами-источниками, 
из которых эти аргументы поступают. 
Затем такие отображения сравниваются для всех значений группы. 
Требуется, чтобы аргументы на одинаковых позициях происходили из одних и тех же групп компоненты. 
Для коммутативных операций порядок аргументов не учитывается. 
Проверка гарантирует, что все линии группы используют одинаковую структуру зависимостей и не требуют дополнительных перестановок операндов.


\subsection{Планирование, легализация и генерация кода}
\label{sec:codegen}

\paragraph{Легализация.}
Для упрощения генерации кода выполняется обычная топологическая сортировка.

\paragraph{Легализация.}
Полученный упорядоченный набор может содержать группы, длина которых превышает 2 --- 
аппаратную ширину 128-битного векторного регистра для 64-битных элементов. 
Процедура разбивает каждую такую группу на подгруппы по 2 элемента. Например, $p = \{v1, v2, v3, v4\}$ 
превращается в две подгруппы $p_1 = \{v1, v2\}$ и $p_2 = \{v3, v4\}$. 
Если длина группы не кратна 2, легализация считается неуспешной, и компонента не векторизуется.

\paragraph{Генерация векторного кода.}
Для каждой группы создаётся одна векторная SSA-инструкция. 
Если аргумент исходной скалярной инструкции принадлежит другой группе, он заменяется на соответствующее векторное значение. 
Если нет --- например, для указателей и памяти --- аргумент передаётся как есть.

После создания всех векторных инструкций выполняется замена скалярных использований:
\begin{itemize}
    \item Внутренние использования в других группах уже учтены на этапе формирования аргументов.
    \item Для сохранений обновляется аргумент \texttt{Memory} у последующих инструкций: 
    вместо исходного значения памяти, порождённого скалярным сохранением, подставляется память, возвращённая векторным сохранением.
    \item Для прочих внешних использований вставляются соответствующие инструкции, извлекающие отдельные элементы из векторного регистра.
\end{itemize}
После этих замен исходные скалярные инструкции становятся «мёртвым кодом» и удаляются.

\subsection{Принятые компромиссы и направления развития}
\label{sec:tradeoffs}

При проектировании был сделан ряд осознанных упрощений, позволивших получить работоспособный векторизатор в условиях ограниченной инфраструктуры Go:

\begin{itemize}
    \item \textbf{Поддерживаемые типы и архитектура.} Только 128-битные векторы.
    Этого достаточно для демонстрации принципиальной работоспособности и ускорения ряда вычислительных шаблонов;
    расширение на 256-битные векторы является технической задачей, не требующей изменения алгоритмического ядра.
    
    \item \textbf{Отсутствие модели стоимости.} Это упрощает код и снижает накладные расходы,
    но может приводить к упущенным оптимизациям. Cитуации, требующие перестановки элементов вектора 
    (например, если одна линия использует аргумент из другой линии), не векторизуются.
    Для первой консервативной реализации можно обойтись без модели стоимости.
    
    \item \textbf{Консервативная модель памяти.} Требование единого \texttt{Memory}-аргумента или 
    строгой цепочки для сохранений исключает многие потенциально векторизуемые случаи, 
    где между двумя независимыми загрузками находится запись по непересекающемуся адресу. 
    Преодоление этого ограничения потребовало бы внедрения более сложного анализа.
    
\end{itemize}

Дальнейшие направления развития включают:
\begin{itemize}
    \item введение упрощённой модели стоимости, оценивающей количество вставок/извлечений, которая позволила бы более агрессивно векторизовывать код;
    \item реализацию базовых \texttt{shuffle}-операций для поддержки некогерентных линий;
    \item улучшение анализа памяти для распознавания независимых загрузок с разными \texttt{Memory}-аргументами;
    \item векторизация редукций.
\end{itemize}

Перечисленные компромиссы и направления формируют реалистичный план развития, сохраняя текущую реализацию компактной и пригодной для включения в основной компилятор Go.